\documentclass[10pt,a4paper]{amsart}
\usepackage{amsmath,amssymb,amsthm,graphicx}
% Retained minimal preamble for the Stacks Project chapter varieties.tex.
% Verbatim excerpt of the chapter preamble: Xy-pic setup (native preamble lines 14-18), the native theorem
% declarations with the mathematical macros (native preamble lines 154-217), and the project own no-op
% definitions of the reference/slogan/history display environments (native wrapper.tex lines 1-3) that the
% counted statement uses for its slogan. Copied mechanically; nothing added, removed or redefined.
% For commutative diagrams we use Xy-pic
\usepackage[all]{xy}

% We use 2cell for 2-commutative diagrams.
\xyoption{2cell}
% Theorem environments.
%
\theoremstyle{plain}
\newtheorem{theorem}[subsection]{Theorem}
\newtheorem{proposition}[subsection]{Proposition}
\newtheorem{lemma}[subsection]{Lemma}

\theoremstyle{definition}
\newtheorem{definition}[subsection]{Definition}
\newtheorem{example}[subsection]{Example}
\newtheorem{exercise}[subsection]{Exercise}
\newtheorem{situation}[subsection]{Situation}

\theoremstyle{remark}
\newtheorem{remark}[subsection]{Remark}
\newtheorem{remarks}[subsection]{Remarks}

\numberwithin{equation}{subsection}

% Macros
%
\def\lim{\mathop{\mathrm{lim}}\nolimits}
\def\colim{\mathop{\mathrm{colim}}\nolimits}
\def\Spec{\mathop{\mathrm{Spec}}}
\def\Hom{\mathop{\mathrm{Hom}}\nolimits}
\def\Ext{\mathop{\mathrm{Ext}}\nolimits}
\def\SheafHom{\mathop{\mathcal{H}\!\mathit{om}}\nolimits}
\def\SheafExt{\mathop{\mathcal{E}\!\mathit{xt}}\nolimits}
\def\Sch{\mathit{Sch}}
\def\Mor{\mathop{\mathrm{Mor}}\nolimits}
\def\Ob{\mathop{\mathrm{Ob}}\nolimits}
\def\Sh{\mathop{\mathit{Sh}}\nolimits}
\def\NL{\mathop{N\!L}\nolimits}
\def\CH{\mathop{\mathrm{CH}}\nolimits}
\def\proetale{{pro\text{-}\acute{e}tale}}
\def\etale{{\acute{e}tale}}
\def\QCoh{\mathit{QCoh}}
\def\Ker{\mathop{\mathrm{Ker}}}
\def\Im{\mathop{\mathrm{Im}}}
\def\Coker{\mathop{\mathrm{Coker}}}
\def\Coim{\mathop{\mathrm{Coim}}}

% Boxtimes
%
\DeclareMathSymbol{\boxtimes}{\mathbin}{AMSa}{"02}

%
% Macros for moduli stacks/spaces
%
\def\QCohstack{\mathcal{QC}\!\mathit{oh}}
\def\Cohstack{\mathcal{C}\!\mathit{oh}}
\def\Spacesstack{\mathcal{S}\!\mathit{paces}}
\def\Quotfunctor{\mathrm{Quot}}
\def\Hilbfunctor{\mathrm{Hilb}}
\def\Curvesstack{\mathcal{C}\!\mathit{urves}}
\def\Polarizedstack{\mathcal{P}\!\mathit{olarized}}
\def\Complexesstack{\mathcal{C}\!\mathit{omplexes}}
% \Pic is the operator that assigns to X its picard group, usage \Pic(X)
% \Picardstack_{X/B} denotes the Picard stack of X over B
% \Picardfunctor_{X/B} denotes the Picard functor of X over B
\def\Pic{\mathop{\mathrm{Pic}}\nolimits}
\def\Picardstack{\mathcal{P}\!\mathit{ic}}
\def\Picardfunctor{\mathrm{Pic}}
\def\Deformationcategory{\mathcal{D}\!\mathit{ef}}
\newenvironment{reference}{}{}
\newenvironment{slogan}{}{}
\newenvironment{history}{}{}
\usepackage{cleveref}
\crefname{originallemma}{original lemma}{original lemmas}
\Crefname{originallemma}{Original lemma}{Original lemmas}
\makeatletter
\newlabel{lemma-closed-fixed-by-Galois}{{038B}{source}{}{sourcepointer-1}{}}
\newlabel{lemma-closed-fixed-by-Galois@cref}{{[originallemma][0][]038B}{[1][0][]source}{}{}{}}
\makeatother
\crefname{originallemma}{original lemma}{original lemmas}
\Crefname{originallemma}{Original lemma}{Original lemmas}
\makeatletter
\newlabel{lemma-galois-action-irreducible-components}{{038J}{source}{}{sourcepointer-2}{}}
\newlabel{lemma-galois-action-irreducible-components@cref}{{[originallemma][0][]038J}{[1][0][]source}{}{}{}}
\makeatother
\crefname{originallemma}{original lemma}{original lemmas}
\Crefname{originallemma}{Original lemma}{Original lemmas}
\makeatletter
\newlabel{lemma-image-irreducible}{{04KX}{source}{}{sourcepointer-3}{}}
\newlabel{lemma-image-irreducible@cref}{{[originallemma][0][]04KX}{[1][0][]source}{}{}{}}
\makeatother
\crefname{originallemma}{original lemma}{original lemmas}
\Crefname{originallemma}{Original lemma}{Original lemmas}
\makeatletter
\newlabel{morphisms-lemma-locally-finite-type-characterize}{{01T2}{source}{}{sourcepointer-4}{}}
\newlabel{morphisms-lemma-locally-finite-type-characterize@cref}{{[originallemma][0][]01T2}{[1][0][]source}{}{}{}}
\makeatother
\title{Galois action on irreducible components of a locally finite type scheme: continuity, finiteness of the defining extension, and orbits}
\author{The Stacks Project Authors}
\date{stacks-project:varieties.tex:local-snapshot-20261011; GNU FDL 1.2 or later}
\begin{document}
\maketitle
\noindent\textit{Editorial note.} All delivered mathematical text is the author's own, copied from the pinned Stacks Project chapter file varieties.tex; no statement or proof step was rewritten, reordered or removed.\par\medskip
\noindent\textit{Editorial note.} The proof is complete as written and ends with the identification W = T; no step is referred to the reader.\par\medskip
\noindent\textit{Editorial note.} The counted spans carry 4 navigation stubs and no bibliography citation.\par\medskip
\noindent\textit{Editorial note.} The preamble is the chapter's retained minimal preamble (sources/stacks-varieties/retained-preamble.tex lines 1-77), a verbatim excerpt of the native preamble.\par\medskip
\noindent\textit{Editorial note.} The label lemma-galois-action-irreducible-components-locally-finite-type is defined only in varieties.tex, so no cross-chapter restatement had to be weighed for this card.\par\medskip
\noindent\textit{Allowed standard prerequisites (editorial).}\begin{itemize}
\item Locally of finite type means affine opens are finite type algebras. Morphisms, Lemma morphisms-lemma-locally-finite-type-characterize, invoked at proof lines 1815-1816. Supplies that the affine chart A is a finite type k-algebra, hence that A tensor the closure is Noetherian.
\item Irreducible components of the base change and their defining ideals. Same chapter, Lemma lemma-galois-action-irreducible-components, invoked at proof lines 1827-1831. Supplies that a component invariant under an open subgroup comes from a component over a finite subextension.
\item Descent of Galois invariant closed subsets. Same chapter, Lemma lemma-closed-fixed-by-Galois, invoked at proof lines 1842-1843. Supplies that a Galois invariant closed subset of the base change is itself a base change of a closed subset of X.
\item Existence of a component of the base change mapping onto a given component. Same chapter, Lemma lemma-image-irreducible, invoked at proof lines 1835-1837. Supplies the component used in part (3).
\end{itemize}\medskip
\ifdefined\hypertarget\hypertarget{sourcepointer-1}{}\fi
\noindent\textit{Source reference.} Original lemma 038B: sources/stacks-varieties/pinned-original-source.tex, source lines 899--899. Pointer outside the counted span, invoked once for context; the load-bearing argument of the counted statement is the author's proof reproduced in full in the proof span.\par\medskip
\ifdefined\hypertarget\hypertarget{sourcepointer-2}{}\fi
\noindent\textit{Source reference.} Original lemma 038J: sources/stacks-varieties/pinned-original-source.tex, source lines 1700--1700. Pointer outside the counted span, invoked once for context; the load-bearing argument of the counted statement is the author's proof reproduced in full in the proof span.\par\medskip
\ifdefined\hypertarget\hypertarget{sourcepointer-3}{}\fi
\noindent\textit{Source reference.} Original lemma 04KX: sources/stacks-varieties/pinned-original-source.tex, source lines 1622--1622. Pointer outside the counted span, invoked once for context; the load-bearing argument of the counted statement is the author's proof reproduced in full in the proof span.\par\medskip
\ifdefined\hypertarget\hypertarget{sourcepointer-4}{}\fi
\noindent\textit{Source reference.} Original lemma 01T2: sources/stacks-chow/tags.txt, source lines 2359--2359. Pointer outside the counted span, invoked once for context; the load-bearing argument of the counted statement is the author's proof reproduced in full in the proof span.\par\medskip
\begin{lemma}
\label{lemma-galois-action-irreducible-components-locally-finite-type}
Let $k$ be a field.
Assume $X \to \Spec(k)$ locally of finite type.
In this case
\begin{enumerate}
\item the action
$$
\text{Gal}(\overline{k}/k)^{opp} \times \text{IrredComp}(X_{\overline{k}})
\longrightarrow
\text{IrredComp}(X_{\overline{k}})
$$
is continuous if we give $\text{IrredComp}(X_{\overline{k}})$ the discrete
topology,
\item every irreducible component of $X_{\overline{k}}$
can be defined over a finite extension of $k$, and
\item given any irreducible component $T \subset X$ the scheme
$T_{\overline{k}}$ is a finite union of irreducible components of
$X_{\overline{k}}$ which are all in the same
$\text{Gal}(\overline{k}/k)$-orbit.
\end{enumerate}
\end{lemma}

\begin{proof}
Let $\overline{T}$ be an irreducible component of $X_{\overline{k}}$.
We may choose an affine open $U \subset X$ such that
$\overline{T} \cap U_{\overline{k}}$ is not empty.
Write $U = \Spec(A)$, so $A$ is a finite type $k$-algebra, see
Morphisms, Lemma \ref{morphisms-lemma-locally-finite-type-characterize}.
Hence $A_{\overline{k}}$ is a finite type $\overline{k}$-algebra,
and in particular Noetherian. Let $\mathfrak p = (f_1, \ldots, f_n)$
be the prime ideal corresponding to $\overline{T} \cap U_{\overline{k}}$.
Since $A_{\overline{k}} = A \otimes_k \overline{k}$
we see that there exists a finite subextension
$\overline{k}/k'/k$ such that each $f_i \in A_{k'}$.
It is clear that $\text{Gal}(\overline{k}/k')$
fixes $\overline{T}$, which proves (1).

\medskip\noindent
Part (2) follows by applying
Lemma \ref{lemma-galois-action-irreducible-components} (1)
to the situation over $k'$ which implies the irreducible component
$\overline{T}$ is of the form $T'_{\overline{k}}$ for some irreducible
$T' \subset X_{k'}$.

\medskip\noindent
To prove (3), let $T \subset X$ be an irreducible component.
Choose an irreducible component $\overline{T} \subset X_{\overline{k}}$
which maps to $T$, see
Lemma \ref{lemma-image-irreducible}.
By the above the orbit of $\overline{T}$ is finite, say it is
$\overline{T}_1, \ldots, \overline{T}_n$. Then
$\overline{T}_1 \cup \ldots \cup \overline{T}_n$
is a $\text{Gal}(\overline{k}/k)$-invariant closed subset of $X_{\overline{k}}$
hence of the form $W_{\overline{k}}$ for some $W \subset X$ closed by
Lemma \ref{lemma-closed-fixed-by-Galois}.
Clearly $W = T$ and we win.
\end{proof}

\end{document}
